% ======================================================================
\section{Discussion and open points}\label{sec:discussion}
% ======================================================================

\subsection{What is new}\label{sec:disc_new}

\begin{itemize}
\item \emph{The operator $\Omega$ to the order required by the NLO cross section,} with the vanishing energy denominators of the dressed one-gluon state defined as principal values, and with the operator ordering of the three-charge part of the coefficient $\mf A^{(3)}_2$, $\mf A^{(1)}_2\mf N_2$, which differs from the order written in~\cite{Lublinsky:2016meo} and on which the cancellations of Section~\ref{sec:rapidity} depend (Section~\ref{sec:omega}).
\item \emph{A complete organisation of the single inclusive spectrum at mid rapidity at order $g^4$} into fifteen products of coefficients, with their longitudinal supports (Section~\ref{sec:groups}).
\item \emph{The fixed-order verification of rapidity factorisation with operator-valued charges:} the two logarithms multiply exactly the BFKL evolution of the projectile correlator and the JIMWLK evolution of the target operator, the latter for an arbitrary ordered correlator, and the terms with more charges cancel in the symmetrised sense (Section~\ref{sec:rapidity}).
\item \emph{The account of all poles,} with the identification of which are ultraviolet and which collinear, and the resulting statements on the scale of the coupling and on the scheme of the fragmentation function (Section~\ref{sec:dglap_rc}).
\item \emph{The finite terms,} row by row, for ordered charges and averaged, with the conditions under which they are integrable (Section~\ref{sec:finite}, Appendices~\ref{app:finite} and~\ref{app:tails}).
\item \emph{The role of the non-commutativity of the charges in the finite terms:} contact terms of the same order as the quadratic terms, and the split of the cubic terms into a part proportional to the C-odd three-point function of the projectile, odd in $\v k$, and a part proportional to the two-point function, even in $\v k$ (Section~\ref{sec:fin_singlet}).
\end{itemize}

\subsection{Structure of the result}\label{sec:disc_lessons}

\paragraph{One longitudinal variable.} The second gluon has a longitudinal momentum $p^+$ between $\Lambda$ and $\vee$. For $p^+\gg k^+$ it belongs to the projectile and gives the BFKL term; for $p^+\ll k^+$ it is indistinguishable from the target and gives the JIMWLK term; for $p^+\sim k^+$ it is either collinear to the measured gluon, which gives the DGLAP term, or a new particle in the final state or in a loop, which gives the finite terms. The three evolutions of the problem are the three boundaries of one phase space, and Row~II.IV alone exhibits all of them (Section~\ref{sec:dg_origin}).

\paragraph{The measured gluon is a dressed gluon.} Counting dressed instead of bare gluons is what makes the spectrum vanish without a target, what brings the coefficients $\mf B_3$ and $\mf C_1$ of the dressed one-gluon state into the amplitudes, and what cancels the ultraviolet part of the self-energy of the measured gluon against the same loop taken across the target. As a consequence the coupling of the emission vertex is renormalised by the wave-function renormalisation of the measured gluon alone.

\paragraph{The charges are operators.} At leading order the commutators of the charges drop out after the average over a colour-neutral projectile, and the picture of classical sources is sufficient. At order $g^4$ they do not: without them the four-charge terms would not reduce to the evolution, the coefficients of the logarithms would not be the BFKL and JIMWLK ones, and the contact terms of the finite part, which the dilute limit requires, would be missing.

\paragraph{The finite terms are defined together with the way in which they are integrated.} Several finite terms are not separately integrable over the position of the unmeasured gluon. The result consists of the terms and of the statements of Section~\ref{sec:fin_together} about which of them stand under a common integral and in which order the integrals are done. These statements follow from the physics (a real and a virtual soft gluon; the measurement as a Fourier integral) and are not additional choices.

\subsection{Use of the result and requirements of a numerical evaluation}\label{sec:disc_use}

\paragraph{Inputs.} The projectile enters through the two-point function $n(\v x,\v y)$ and, for the terms odd in $\v k$, through the symmetrised three-point function $\sigma(\v x,\v y,\v z)$; for a projectile that is not globally colour invariant the ordered forms of Appendix~\ref{app:finite} have to be used with the full ordered correlators. The target enters through averages of traces of adjoint Wilson lines. At leading order these are adjoint dipoles. The JIMWLK term and the finite terms require averages of products of more than two Wilson lines at different transverse positions (three in the symmetrised cubic terms; up to four in some two-charge terms as they are written), which are not determined by the dipole alone; the C-odd traces with $D^h$ require a target ensemble that is not C-even to contribute.

\paragraph{Evolution.} The logarithms $\log(\vee/k^+)$ and $\log(k^+/\Lambda)$ are to be absorbed into the evolution of the projectile correlator and of the target averages up to the rapidity of the measured gluon (Section~\ref{sec:logsummary}). A description of NLO accuracy throughout would require, in addition, the NLO evolution of both; this is outside the scope of this paper. The finite terms refer to cut-offs in the plus momentum; with another definition of the rapidity cut-off they would differ by non-logarithmic terms.

\paragraph{Scales.} The renormalisation and factorisation scales enter through the running-coupling and DGLAP terms only, Eq.~\eqref{eq:dg_muRmuF}; the local form~\eqref{eq:dg_local} is the convenient one in momentum space, and it has no large-distance tail of its own. The spectrum of hadrons is obtained by convolution with $\MS$ fragmentation functions defined as number densities; with fragmentation functions in the convention of~\cite{Chirilli:2012jd} the term $\Delta$ of Eq.~\eqref{eq:dg_Delta} has to be subtracted.

\paragraph{Integrals.} A finite term is an integral over one momentum fraction and up to six transverse positions (for instance Eq.~\eqref{eq:S-III.V-3}), two of which carry the measured momentum. A numerical evaluation has to (i) perform the Fourier integrals at the measured momentum first, or in a way that is equivalent to this; (ii) combine the members of the pairs A--D, and the three rows of the two-gluon final state, before the integral over the position of the unmeasured gluon; (iii) treat the principal values of Row~I.III and the distributions of the instantaneous terms as defined in Section~\ref{sec:nuance_pv}; (iv) use a rotation-invariant regularisation at the collinear point of Row~I.V. A momentum-space form of all finite terms, which would make step (i) trivial, exists at present only in the dilute limit. We have not carried out such an evaluation. No numerical evaluation of the finite terms, and no plot, is given in this paper, and we make no statement about the numerical size, or the sign, of the NLO correction.

\subsection{Limitations and open points}\label{sec:openpoints}

The result is subject to the following restrictions, which are assumptions of the framework and not approximations made along the way: the soft gluons are treated in the eikonal approximation, i.e.\ at leading power in $k^+/\vee$ and $\Lambda/k^+$, with $\vee\to\infty$ taken before $\epsilon\to0$ (Section~\ref{sec:dg_poles}); the projectile is dilute, i.e.\ the result is a fixed-order expansion in the charge density of the projectile, and it does not describe the collision of two dense objects; the theory is pure Yang--Mills theory; the gluon is produced at mid rapidity, far from both the projectile and the target. Within this framework the following points are not settled by this paper.
\begin{enumerate}
\item[(i)] \emph{Dilute limit of the virtual correction.} The comparison of the virtual correction with the one-loop central-emission vertex~\cite{Fadin:1998py,DelDuca:1998cx,Bartels:2006hg} has not been completed. It requires virtual and real parts with matched regulators; in particular, it is open whether the target-side term $\frac{\alpha_sN_c}{\pi}[\frac2\epsilon\log\frac{k^+}{\Lambda}+\log^2\frac{k^+}{\Lambda}]$ of the virtual correction for a point-like quark, Eq.~\eqref{eq:dilute_groupI}, cancels against the real emission from the region $|\v p|\lesssim(\Lambda/k^+)|\v k|$, which we have not computed, and whether the finite function $C(\v k,\v l)$ agrees with the one-loop vertex up to terms that can be absorbed into the unintegrated gluon distributions. Also not done are the conversion of the virtual trajectories between our plus-momentum cut-offs and the energy scales of~\cite{Bartels:2006hg}, the reduction of the virtual correction for a general projectile, the dilute limit of the coordinate-space finite terms as printed, and the comparison of the printed colour-singlet contact terms with the commutator contributions required in this limit.
\item[(ii)] \emph{Comparison with Li and Skokov.} Our symmetrised cubic terms and the first saturation correction of~\cite{Li:2021ntm} agree in their colour structures but not in their parity in $\v k$: their result has a $\v k$-even part, which our symmetrised terms do not have, and in the two-Wilson-line sector their even term has no counterpart here, while our odd two-line term has none there. We traced the difference to a relative phase $i$ in the two-charge structures that build the non-linear Weizs\"acker--Williams field and in the rescattering on the second charge after the collision; its origin, possibly the prescription for $1/k^+$ (retarded or symmetric $x^-$ profile), has not been determined, and the even part of~\cite{Li:2021ntm} has not been checked independently. A kernel-level comparison of the odd three-line terms, and the question whether our odd two-line terms add up to zero, are open.
\item[(iii)] \emph{Prescription for the vanishing energy denominators.} The finite terms are given for the principal-value definition~\eqref{eq:PV} of the vanishing energy denominators of the dressed measured gluon. An adiabatic $i\varepsilon$ prescription, as in~\cite{Radhakrishnan:2026soft}, differs by terms on the energy shell, $\propto\delta(E_i-E_n)$, which would add a finite term to Row~I.III. This term is not included, and the dependence of the result on the prescription has not been studied beyond the dilute limit for a point-like projectile, where it is absent (Section~\ref{sec:nuance_pv}).
\item[(iv)] \emph{Numerical evaluation.} No numerical evaluation of the finite terms is given. Apart from the dilute limit, the finite terms are given in coordinate space only; a momentum-space form of all finite terms, and a complete numerical integration over all transverse positions with realistic Wilson lines, have not been carried out.
\item[(v)] \emph{Constant and $\beta_0$ terms.} The LO-proportional constant that is left in the finite terms of Row~I.V in our grouping has not been evaluated for a general projectile and target, and the momentum-weighted collinear regions of three rows (II.I, III.III, I.V) have not been integrated. The logarithm of the emitter distance times the inverse saturation scale, with a $\beta_0$-type coefficient, that the resolving region of Row~I.V contains by power counting, has not been evaluated with a target model. The $N_f$ terms are not part of this pure Yang--Mills calculation; the quark loop is used only for the identification of the $\beta_0$ terms (Appendix~\ref{app:scheme}).
\item[(vi)] \emph{Collinear limit.} The collinear convention is fixed by the general $\MS$ counterterm and by two angle-integrated coefficient functions; a comparison with the NLO coefficient of hadroproduction, $pp\to gX$, and a collinear limit of the hybrid formula of~\cite{Chirilli:2012jd} have not been carried out. At the conditionally convergent collinear point of Row~I.V, a continuation of the transverse indices to $d_\perp=2-\epsilon$ could give a finite local term; whether it is present has not been decided.
\item[(vii)] \emph{Status of the integrability statements.} The power counting of the contact kernels is numerical (exponent fits at generic configurations, two smooth target fields) together with analytic arguments, not a proof for every degenerate configuration; the kernel of the $f^{abc}\{\rho^b,\rho^c\}$ term of the vertex correction was estimated analytically only. The absence of a contribution from the region in which only part of the gluon pair is far was verified numerically for one charge and a target uniform at large distance; for a colour-neutral projectile and a non-uniform target it rests on an argument. The interchange of the $p^+$ integral with the joint large-distance limit was checked numerically for one of the two rows concerned. The statements on the instantaneous terms were checked with model functions of the correct short-distance behaviour, and the universality of the shift of the BFKL logarithm under a change of the separation of soft and valence modes is argued from the form of the integrands, not checked row by row.
\item[(viii)] \emph{Coefficients quoted or not re-derived.} The coefficients of the dressed vacuum, $\mf B_2$ and $\mf A^{(3)}_2$, are quoted from~\cite{Lublinsky:2016meo}; of the latter only the order of the charges in the three-charge term was determined here. The transformation of $\mf B_2$ and $\mf C_1$ to coordinate space is quoted without derivation (Appendix~\ref{app:coefficients}). The support of the commutator part of the two-emission term of $\mf B_2$ is taken as $k^++p^+<\vee$ (Section~\ref{sec:vacin}).
\end{enumerate}
